\documentclass[12pt, a4paper]{ctexart}
\usepackage{amsmath, amssymb, amsthm}
\usepackage{geometry}
\usepackage{fancyhdr}
\usepackage{enumitem}
\usepackage{graphicx}
\usepackage{tikz}
\usepackage{array}
\usepackage{multirow}
\usepackage{booktabs}

\geometry{left=2cm, right=2cm, top=2.5cm, bottom=2.5cm}
\setCJKmainfont{Hiragino Sans GB}
\setCJKsansfont{Heiti SC}

\pagestyle{fancy}
\fancyhf{}
\fancyhead[L]{培优专题·初等数论}
\fancyhead[R]{测试卷}
\fancyfoot[C]{\thepage}
\renewcommand{\headrulewidth}{0.4pt}

\newcommand{\score}[1]{\hfill\makebox[2cm][l]{（#1 分）}}

\begin{document}

\begin{center}
{\LARGE\bfseries 培优专题·初等数论 测试卷}\\[0.5em]
{\large 总分：150分 \quad 时间：120分钟}
\end{center}

\vspace{1em}

\noindent\textbf{一、单项选择题（本大题共8小题，每小题5分，共40分。每小题只有一个选项符合题意。）}

\vspace{0.5em}

\begin{enumerate}
\item $2^{100}$ 除以 $7$ 的余数为 \score{5}
\begin{enumerate}[label=\Alph*.]
\item $1$
\item $2$
\item $4$
\item $6$
\end{enumerate}

\item 满足 $\gcd(a,b)=6$，$\mathrm{lcm}(a,b)=180$ 的正整数对 $(a,b)$（不计顺序）的个数为 \score{5}
\begin{enumerate}[label=\Alph*.]
\item $2$
\item $3$
\item $4$
\item $6$
\end{enumerate}

\item $100!$ 中因子 $5$ 的最高幂次为 \score{5}
\begin{enumerate}[label=\Alph*.]
\item $20$
\item $24$
\item $25$
\item $30$
\end{enumerate}

\item 用辗转相除法求 $\gcd(2024, 2023)$ 的值为 \score{5}
\begin{enumerate}[label=\Alph*.]
\item $1$
\item $7$
\item $11$
\item $13$
\end{enumerate}

\item 模 $11$ 的二次剩余共有 \score{5} 个。
\begin{enumerate}[label=\Alph*.]
\item $4$
\item $5$
\item $6$
\item $10$
\end{enumerate}

\item 不定方程 $3x+5y=100$ 的正整数解 $(x,y)$ 的个数为 \score{5}
\begin{enumerate}[label=\Alph*.]
\item $5$
\item $6$
\item $7$
\item $8$
\end{enumerate}

\item $2^{100}$ 的十进制表示的末两位数字为 \score{5}
\begin{enumerate}[label=\Alph*.]
\item $76$
\item $24$
\item $56$
\item $44$
\end{enumerate}

\item 设 $p$ 为奇素数，则同余方程 $x^2\equiv 1\pmod{p}$ 的解的个数为 \score{5}
\begin{enumerate}[label=\Alph*.]
\item $1$
\item $2$
\item $p-1$
\item 取决于 $p$ 的值
\end{enumerate}
\end{enumerate}

\vspace{1em}

\noindent\textbf{二、多项选择题（本大题共4小题，每小题5分，共20分。每小题有多个选项符合题意，全部选对得5分，部分选对得2分，有选错得0分。）}

\vspace{0.5em}

\begin{enumerate}[resume]
\item 关于最大公因数和最小公倍数，下列说法正确的是 \score{5}
\begin{enumerate}[label=\Alph*.]
\item $\gcd(a,b)\cdot\mathrm{lcm}(a,b)=|ab|$ 对任意正整数 $a,b$ 成立
\item 若 $\gcd(a,b)=1$，则称 $a,b$ 互素
\item 若 $p$ 为素数且 $p\mid ab$，则 $p\mid a$ 或 $p\mid b$
\item $\gcd(a,b)=\gcd(a,b-a)$ 对任意正整数 $a<b$ 成立
\end{enumerate}

\item 关于同余的性质，下列正确的是 \score{5}
\begin{enumerate}[label=\Alph*.]
\item 若 $a\equiv b\pmod{m}$ 且 $c\equiv d\pmod{m}$，则 $a+c\equiv b+d\pmod{m}$
\item 若 $a\equiv b\pmod{m}$，则 $a^k\equiv b^k\pmod{m}$ 对所有正整数 $k$ 成立
\item 若 $ac\equiv bc\pmod{m}$，则 $a\equiv b\pmod{m}$
\item 欧拉定理：若 $\gcd(a,m)=1$，则 $a^{\varphi(m)}\equiv 1\pmod{m}$
\end{enumerate}

\item 关于费马小定理，以下说法正确的是 \score{5}
\begin{enumerate}[label=\Alph*.]
\item 若 $p$ 为素数且 $p\nmid a$，则 $a^{p-1}\equiv 1\pmod{p}$
\item 费马小定理是欧拉定理的特殊情况
\item 费马小定理可以用于判断素数
\item $a^p\equiv a\pmod{p}$ 对所有整数 $a$ 和素数 $p$ 成立
\end{enumerate}

\item 关于中国剩余定理（孙子定理），下列正确的是 \score{5}
\begin{enumerate}[label=\Alph*.]
\item 当模数两两互素时，同余方程组有唯一解（模各模数之积）
\item $x\equiv 2\pmod{3}$，$x\equiv 3\pmod{5}$，$x\equiv 2\pmod{7}$ 有解
\item 中国剩余定理可以用于大数分解
\item 当模数不两两互素时，方程组一定无解
\end{enumerate}
\end{enumerate}

\vspace{1em}

\noindent\textbf{三、填空题（本大题共4小题，每小题5分，共20分。）}

\vspace{0.5em}

\begin{enumerate}[resume]
\item $\gcd(168, 252)=$ \underline{\hspace{2cm}}，$\mathrm{lcm}(168, 252)=$ \underline{\hspace{2cm}}。\score{5}

\item $7^{200}$ 的末位数字为 \underline{\hspace{2cm}}。\score{5}

\item 中国剩余定理求解：$x\equiv 1\pmod{3}$，$x\equiv 2\pmod{5}$，$x\equiv 3\pmod{7}$，则 $x\equiv$ \underline{\hspace{3cm}} $\pmod{105}$。\score{5}

\item 欧拉函数 $\varphi(360)=$ \underline{\hspace{2cm}}。\score{5}
\end{enumerate}

\vspace{1em}

\noindent\textbf{四、解答题（本大题共6小题，共70分。解答应写出文字说明、证明过程或演算步骤。）}

\vspace{0.5em}

\begin{enumerate}[resume]
\item （10分）求 $2^{2024}+3^{2024}$ 除以 $5$ 的余数。\score{10}

\item （12分）用辗转相除法求 $\gcd(2024, 2023)$，并将其表示为 $2024$ 和 $2023$ 的整系数线性组合，即求整数 $x,y$ 使得 $2024x+2023y=\gcd(2024,2023)$。\score{12}

\item （12分）证明：对任意正整数 $n$，$n^5-n$ 能被 $30$ 整除。\score{12}

\item （12分）求不定方程 $7x+11y=200$ 的所有正整数解。\score{12}

\item （12分）利用欧拉定理或费马小定理求 $7^{222}$ 除以 $11$ 的余数。\score{12}

\item （12分）证明威尔逊定理：$p$ 为素数当且仅当 $(p-1)!\equiv -1\pmod{p}$。\score{12}

\end{enumerate}

\end{document}
